% !TEX program = lualatex
\documentclass[a4paper,11pt]{scrartcl}

\usepackage[
  lang=en,
  mode=journal,
  theme=ocean,
  toc=true,
  toc-layout=page
]{synaptic}
\usepackage{booktabs,tabularx}
\newcolumntype{Y}{>{\raggedright\arraybackslash}X}
\newcommand{\DocCmd}[1]{\textcolor{syn-primary!80!black}{\texttt{\textbackslash#1}}}

\SynapticNewTheorem[remark]{synobservation}{Observation}
\SynapticHeader{Synaptic User Manual}
\SynapticSubHeader{Reliable academic typesetting with LuaLaTeX}
\SynapticPreprint{v5.2.0}
\SynapticTitle{The synaptic Package: User Manual}
\SynapticAuthor{synaptic project}
\SynapticAffiliation{Open-source LaTeX community}
\SynapticEmail{synaptic@example.com}
\SynapticSubject{User manual for synaptic version 5.2.0}
\SynapticKeywords{synaptic, LaTeX3, LuaLaTeX, academic typesetting}

\begin{synabstract}
  \textbf{synaptic} is a modular LaTeX3 design system for journal articles,
  books, lecture material, compact notes, and Beamer presentations. This manual describes its
  installation, requirements, safe configuration model, typographic system,
  metadata API, theorem system, theme and customization tools, and
  mode-specific features.
\end{synabstract}

\begin{document}
\SynapticMakeTitle

\section{Getting started}

The package requires LuaLaTeX. Print-oriented modes use a KOMA-Script class;
presentation mode uses the Beamer class. If the package is installed in your
TeX tree, a journal document needs nothing more than:

\begin{verbatim}
\documentclass{scrartcl}
\usepackage[mode=journal]{synaptic}
\SynapticTitle{A Short Article}
\SynapticAuthor{Ada Author}
\begin{synabstract}Summary.\end{synabstract}
\begin{document}
\SynapticMakeTitle
\section{Introduction}
Text.
\end{document}
\end{verbatim}

Use \texttt{scrbook} for book mode. The selected class controls the paper size;
synaptic does not silently force A4.

For a widescreen presentation, select both the Beamer class and Synaptic's
matching mode:

\begin{verbatim}
\documentclass[aspectratio=169]{beamer}
\usepackage[mode=beamer,theme=midnight,numbering=none]{synaptic}
\SynapticTitle{A Short Presentation}
\SynapticInstructor{Ada Lecturer}
\begin{document}
\SynapticMakeTitle
\section{First idea}
\SynapticSectionFrame
\begin{frame}{A focused frame}Content.\end{frame}
\end{document}
\end{verbatim}

\subsection{Using a checkout without installation}

Extract the package modules, then point \texttt{TEXINPUTS} at the build output:

\begin{verbatim}
l3build unpack
export TEXINPUTS=/path/to/synaptic/build/unpacked//:$TEXINPUTS
lualatex document.tex
\end{verbatim}

\subsection{System installation}

A TeX Live installation can install (and, when versions change, replace) the
bundle from the repository root:

\begin{verbatim}
l3build install   # copy into the user's texmf tree
l3build uninstall # remove a previous installation
l3build ctan      # produce TDS/CTAN archives
\end{verbatim}

Any standard \texttt{texmf} tree works; the generated files belong in
\texttt{tex/latex/synaptic/}. Archives from \texttt{l3build ctan} unpack
directly into the same structure.

\subsection{Requirements}

\begin{itemize}
  \item LuaLaTeX; pdfLaTeX and XeLaTeX are rejected with an explicit error.
  \item A KOMA-Script class (\texttt{scrartcl}, \texttt{scrreprt}, or
    \texttt{scrbook}) for the four print modes, or \texttt{beamer} for
    \texttt{mode=beamer}.
  \item A recent TeX distribution with \texttt{expl3}, \texttt{fontspec},
    \texttt{unicode-math}, \texttt{geometry}, \texttt{microtype},
    \texttt{thmtools}, and the other standard dependencies.
  \item Libertinus Serif, Sans, Mono, and Math as system fonts.
  \item \texttt{tcolorbox} and PGF for teaching and notes cards. For
    \texttt{lang=zh}, install \texttt{ctex}, Source Han Serif/Sans CN, and
    LXGW WenKai.
\end{itemize}

\section{Configuration}

\noindent{\setlength{\parindent}{0pt}\noindent\begin{tabularx}{\textwidth}{@{}l l Y@{}}
  \toprule
  \textbf{Key} & \textbf{Default} & \textbf{Values and purpose} \\
  \midrule
  \texttt{mode} & journal & journal, book, lecture, notes, or beamer \\
  \texttt{theme} & ocean & ocean, graphite, forest, midnight, or paper \\
  \texttt{background} & white & any xcolor expression for the page colour \\
  \texttt{lang} & en & en or zh; selects labels and CJK support \\
  \texttt{toc} & false & whether the title workflow prints a TOC \\
  \texttt{toc-layout} & top & top, inline, or page \\
  \texttt{title-layout} & auto & inline, compact, page, sequence, or mode default \\
  \texttt{numbering} & auto & section, chapter, continuous, none, or mode default \\
  \texttt{color-mode} & screen & screen, print, or mono \\
  \texttt{density} & balanced & compact, balanced, or airy \\
  \texttt{measure} & auto & narrow, standard, wide, or mode default \\
  \texttt{theorem-style} & boxed & boxed or plain \\
  \texttt{binding-offset} & 0pt & additional book binding correction \\
  \bottomrule
\end{tabularx}}

Only \texttt{theme}, \texttt{background}, \texttt{toc}, and \texttt{toc-layout}
are safe runtime settings:

\begin{verbatim}
\SynapticSetup{theme=paper,background=gray!20,toc=false}
\end{verbatim}

Mode, language, and load-time typography options are frozen after package
loading. Attempts to change them through \texttt{\string\SynapticSetup} raise an
error instead of producing a half-reconfigured document.

\subsection{Modes}

\noindent{\setlength{\parindent}{0pt}\noindent\begin{tabularx}{\textwidth}{@{}l l Y@{}}
  \toprule
  \textbf{Mode} & \textbf{Class} & \textbf{Purpose} \\
  \midrule
  journal & scrartcl & Inline article title, metadata, running heads, statements \\
  book & scrbook & Book title sequence, parts, chapters, duplex navigation \\
  lecture & scrartcl/scrreprt & Course masthead and semantic teaching system \\
  notes & scrartcl & Compact masthead and lightweight knowledge cards \\
  beamer & beamer & Projection-first frames, agenda/section slides, progress navigation \\
  \bottomrule
\end{tabularx}}

\textbf{journal} renders an inline article heading on the first page. The
running head carries the first author and short title, and the block supports
abstract, keywords, arXiv, correspondence, publication information, and
statement sections for journal submissions.

\textbf{book} renders a half title, formal title page, verso block, and
optional dedication, then provides front/main/back matter helpers, part pages,
chapter subtitles, epigraphs, and mirrored duplex running heads.

\textbf{lecture} renders a course masthead (course code, title, lecturer,
semester, lecture number) and wraps theorems and examples in a sharp card
system. Learning goals and lecture summaries frame each teaching unit.

\textbf{notes} renders a compact masthead with an optional tag line and
supports class-level two-column composition with balanced final columns.

\textbf{beamer} renders a projection-first title frame, semantic frame titles,
single-layer semantic cards, and a restrained footer with frame progress. It
reuses course metadata and the lecture teaching-card vocabulary. Explicit
\DocCmd{SynapticSectionFrame} dividers keep section pacing under the author's
control instead of inserting slides implicitly. The standard
\DocCmd{titlegraphic} typesets a graphic right-aligned in the bottom row of the
title frame, on the tag line when tags are present. Namespaced statement
environments render as one tcolorbox card each, with the amsthm-style head
line typeset inside the box, so a Beamer block is never nested in the card
(which previously reproduced the card chrome and added height). Native
\texttt{theorem}, \texttt{definition}, and \texttt{example} environments keep
Beamer's default block template.

\subsection{Fonts}

Synaptic deliberately uses one stack on every machine: Libertinus Serif,
Sans, Mono, and Math for Latin text and mathematics. Because Libertinus Math
has no separate bold face, bold mathematics uses a deterministic synthetic
weight of the same design. Source Han Serif/Sans CN
for Simplified Chinese; and LXGW WenKai for emphasized Chinese text. Missing
fonts are configuration errors, so typography never changes silently.

\section{The typographic system}

Synaptic's default appearance is deterministic rather than decorative. A
handful of design decisions repeat throughout every mode.

\subsection{Colour delivery}

{\setlength{\parindent}{0pt}\noindent\begin{tabularx}{\textwidth}{@{}l l Y@{}}
  \toprule
  \textbf{color-mode} & \textbf{Use} & \textbf{Effect} \\
  \midrule
  screen & on-screen reading & full theme colours, coloured links \\
  print & paper output & theme colours, dark (print-safe) links \\
  mono & monochrome delivery & black-grey tokens only; semantic colours map
  to greys \\
  \bottomrule
\end{tabularx}}

\subsection{Page background}

\textbf{background} selects the page colour with any xcolor expression and is
runtime-adjustable:

\begin{verbatim}
\usepackage[background=blue!10]{synaptic}
\SynapticSetup{background=midnight!85!black}
\end{verbatim}

Every derived role knows the active background. Cards, borders, hairlines, and
title rules blend against the literal page colour, so a light tinted document
stays coherent without touching the theme. When the relative luminance of the
background falls below the mid-point, body and muted text flip to light inks,
rules and borders are lightened instead of darkened, and links use a lighter
tint of the primary colour. In Beamer mode the frame background canvas follows
the same setting, so dark presentations are supported without extra setup.
Restore the default with \texttt{background=white}.

\subsection{Density and measure}

\textbf{density} tunes line spacing in three steps: \texttt{compact} for dense
reference material, \texttt{balanced} as the mode-aware default, and
\texttt{airy} for slide-like handouts. \textbf{measure} overrides the
mode-default column width: \texttt{narrow} for long prose, \texttt{standard}
for journals, \texttt{wide} for design-driven layouts. In book mode, measure
adjusts the mirrored inner/outer margins instead of breaking duplex
symmetry.

\subsection{Title composition}

\textbf{title-layout} selects how each mode renders its title area. The
\texttt{auto} value resolves to \texttt{inline} (journal), \texttt{sequence}
(book), \texttt{compact} (lecture, notes), and \texttt{page} (beamer). \texttt{page} dedicates a
full title page and \texttt{inline} forces the first-page composition.
Beamer accepts \texttt{auto}, \texttt{page}, and \texttt{compact}; the last
reduces the title scale for metadata-heavy decks.
All title renderers reuse an asymmetric “synaptic rail”: a short primary rule,
an accent node, and a quiet continuation line. Book title and verso pages use
the motif at architectural edges; lecture and notes place it beside compact
metadata panels, while journal mode keeps the treatment editorial and inline.

\subsection{Navigation}

Every paged mode draws a hairline beneath its running heads. The hairline uses a
separate, deliberately lighter decorative token so navigation stays quiet.
Running-head labels and outer folios repeat the small node marker; book mode
additionally gives forced duplex blank pages a quiet footer anchor.
Beamer echoes the same hairline beneath each frame title and uses a clamped
progress rule below its footer, so the first compilation is safe before the
final frame count is known.

\subsection{Headings, contents, and captions}

All numbered section levels share the sans heading family:
\texttt{section} uses the primary colour, \texttt{subsection} the dark text
colour, and \texttt{subsubsection} a
smaller, slightly muted tint, so structure reads at a glance. When a table of
contents is requested its top-level entries repeat the on-page heading colour
in bold while deeper entries stay dark, and the number alignment and vertical
rhythm are tightened. Figure and table caption labels carry the secondary
theme token to keep them consistent with the rest of the design.

\subsection{Cards and rules}

Teaching and notes cards share one visual language: a very small corner radius,
a west accent bar tinted with the semantic colour of the card (primary,
secondary, warning, or success), a faint hairline, and a light surface with a
readable title colouring.
Rules and hairlines derive from the theme's primary colour at low tint, so
decoration never competes with text.

\section{Title and PDF metadata}

All five modes share the metadata model and render it differently:

\noindent{\setlength{\parindent}{0pt}\noindent\begin{tabularx}{\textwidth}{@{}l Y@{}}
  \toprule
  \textbf{Command} & \textbf{Meaning} \\
  \midrule
  \DocCmd{SynapticTitle} & Required printed title and PDF title \\
  \DocCmd{SynapticSubtitle} & Optional display subtitle \\
  \DocCmd{SynapticShortTitle} & Running-head title \\
  \DocCmd{SynapticDate} & Explicit date \\
  \DocCmd{SynapticAuthor} & Repeatable author and optional affiliation mark \\
  \DocCmd{SynapticAffiliation} & Repeatable affiliation \\
  \DocCmd{SynapticEmail} & Repeatable clickable contact address \\
  \DocCmd{SynapticSubject} & PDF subject \\
  \DocCmd{SynapticKeywords} & Printed and PDF keywords \\
  \DocCmd{SynapticHeader} & Submission header override \\
  \DocCmd{SynapticSubHeader} & Secondary header \\
  \DocCmd{SynapticPreprint} & Preprint identifier \\
  \DocCmd{SynapticArxiv} & Linked arXiv identifier \\
  \DocCmd{SynapticDedicated} & Dedication \\
  \DocCmd{SynapticMakeTitle} & Validate and render all title data \\
  \bottomrule
\end{tabularx}}

Journal mode accepts publication information, correspondence, and author notes;
book mode accepts half-title, publisher, edition, ISBN, and copyright fields;
lecture mode accepts course, code, lecture number, instructor, and semester.
Beamer consumes the same course fields and synchronizes title, author,
subject, and keywords with Beamer's native PDF metadata path.
Use \verb|\SynapticStatement{heading}{body}| for journal funding, data, ethics,
or conflict statements; the starred form omits its table-of-contents entry.
PDF metadata receives one Unicode-safe conversion.

\section{Statements and mathematics}

The built-in \texttt{syntheorem}, \texttt{synlemma}, \texttt{synproposition},
\texttt{syndefinition}, \texttt{syncorollary}, \texttt{synaxiom}, \texttt{synconjecture},
\texttt{synclaim}, \texttt{synproperty}, \texttt{syncriterion}, \texttt{synconvention},
\texttt{synproblem}, \texttt{synsolution}, \texttt{synexample}, \texttt{synexercise},
\texttt{synremark}, \texttt{synnote}, and \texttt{synwarning} environments share one
counter. Books use chapter scope by default; shorter modes use sections.
Definitions use upright body text; proposition-style statements use italics in
English and upright CJK text. Examples and exercises are plain-style
statements; warnings use the remark presentation.

By default (\texttt{theorem-style=boxed}) the numbered statement environments
are presented as refined cards: a subtle theme-tinted background, a semantic
west accent bar, a faint hairline, and the same restrained corner treatment as
the teaching and notes card language.
\texttt{synproof} stays clean and unboxed. The standard \texttt{proof}
environment remains untouched. Select \texttt{theorem-style=plain} for
the classic unboxed look.

\begin{syntheorem}[Gaussian integral]
  The integral satisfies
  \[
    \int_0^\infty e^{-x^2}\,\mathrm{d}x=\frac{\sqrt{\pi}}{2}.
  \]
\end{syntheorem}

\begin{synproof}
  Square the integral, change to polar coordinates, and evaluate the radial
  integral.
\end{synproof}

New theorem types accept a style:

\begin{verbatim}
\SynapticNewTheorem[remark]{synobservation}{Observation}
\end{verbatim}

The available styles are \texttt{plain}, \texttt{definition}, and
\texttt{remark}; a custom statement is boxed automatically. Its environment
name must be a lowercase alphanumeric identifier beginning with \texttt{syn}.
Duplicate or unnamespaced names are errors.

\begin{synobservation}
  Unicode-math commands such as \(\symbb{R}\), \(\symcal{F}\), and
  \(\symfrak{g}\) are available directly, without generated letter aliases.
\end{synobservation}

\section{Themes and customization}

Themes are switched through \verb|\SynapticSetup{theme=name}|. Define an
immutable custom theme with xcolor expressions, then activate it explicitly:

\begin{verbatim}
\SynapticDefineTheme{brand}{blue!70!black}{gray!80!black}[orange]
\SynapticSetup{theme=brand}
\end{verbatim}

The arguments are a lowercase slug, primary and secondary colours, and an
optional accent. Existing theme names cannot be replaced.
The semantic tokens below are always available for extensions and always
derive from the active theme:

{\setlength{\parindent}{0pt}\noindent\begin{tabularx}{\textwidth}{@{}l l Y@{}}
  \toprule
  \textbf{Token} & \textbf{Role} & \textbf{Use} \\
  \midrule
  \texttt{syn-primary} & brand colour & headings, title rules, accents \\
  \texttt{syn-secondary} & ink colour & muted text, remark tones \\
  \texttt{syn-accent} & accent colour & optional emphasis \\
  \texttt{syn-text-dark} & body ink & readable text \\
  \texttt{syn-muted-text} & readable muted & auxiliary labels \\
  \texttt{syn-surface} & card background & boxes and surfaces \\
  \texttt{syn-border} & card border & box frames \\
  \texttt{syn-rule} & decorative hairline & separators \\
  \texttt{syn-bg} & page background & follows the \texttt{background} key \\
  \texttt{syn-danger} & semantic danger & warnings on the red side \\
  \texttt{syn-warning} & semantic warning & warnings on the amber side \\
  \texttt{syn-success} & semantic success & exercises, todo lists \\
  \bottomrule
\end{tabularx}}

Lecture and Beamer modes include two teaching cards:
\begin{itemize}
  \item \texttt{synlearninggoals} for the goals of a session;
  \item \texttt{synlecturesummary} for its closing summary.
\end{itemize}
The course metadata commands described above complete the lecture workflow.
Notes mode provides \texttt{synkeyidea}, \texttt{synquestion},
\texttt{syntodo}, and \texttt{synsummary}.

\subsection{Beamer agenda and section workflow}

With \texttt{toc=true}, \texttt{toc-layout=inline} places a compact agenda on
the title frame, \texttt{top} creates a regular numbered agenda frame, and
\texttt{page} creates a plain unnumbered agenda frame. Use
\DocCmd{SynapticAgendaFrame}\texttt{[Optional title]} to insert another agenda
later. After a normal \verb|\section{...}|, call
\DocCmd{SynapticSectionFrame} when a visual divider is useful. Plain title,
agenda, and divider frames do not inflate the visible progress count.

\section{Book workflow and acknowledgments}

Use \DocCmd{SynapticFrontMatter}, \DocCmd{SynapticMainMatter}, and
\DocCmd{SynapticBackMatter} around the corresponding parts of a book.
\DocCmd{SynapticBookContents}, chapter subtitles, epigraphs, and
list/appendix helpers complete the long-form workflow. Notes mode provides
\texttt{synkeyidea}, \texttt{synquestion}, \texttt{syntodo}, and
\texttt{synsummary}; class-level two-column layouts are supported. The
\DocCmd{SynapticAcknowledgments} command creates a chapter in book mode and
a section in other modes. Its starred form suppresses the TOC entry.

\section{Troubleshooting}

\begin{itemize}
  \item \textbf{``The synaptic package requires LuaLaTeX''} --- you are running
    pdfLaTeX or XeLaTeX. Engines other than LuaLaTeX are rejected on purpose.
  \item \textbf{``requires a KOMA-Script document class''} --- use a KOMA class
    with \texttt{scrartcl}, \texttt{scrreprt}, or \texttt{scrbook}.
  \item \textbf{Beamer class mismatch} --- use the \texttt{beamer} document
    class whenever \texttt{mode=beamer} is selected.
  \item \textbf{``mode=book requires a chapter-based class''} --- use
    \texttt{scrbook} (or another KOMA class defining \texttt{chapter} and
    \texttt{frontmatter}).
  \item \textbf{``Required font ... not found''} --- install the documented
    system font packages; Synaptic does not substitute a different design.
  \item \textbf{``Options frozen after loading''} --- \texttt{mode},
    \texttt{lang}, and typography keys are structural.
    Re-select them in \verb|\usepackage|; runtime changes via
    \verb|\SynapticSetup| are errors.
  \item \textbf{No CJK font found} --- install Source Han Serif/Sans CN and
    LXGW WenKai.
  \item \textbf{``Document title is empty''} --- set
    \verb|\SynapticTitle| before calling \verb|\SynapticMakeTitle|.
  \item \textbf{Paper size unwanted} --- the package never overrides the class
    paper size; select it in the class options.
\end{itemize}

\section{Development commands}

\begin{verbatim}
./scripts/verify # tests, manuals, examples, diagnostic scan
l3build unpack   # extract .sty and .def files from synaptic.dtx
l3build ctan     # create tested TDS and CTAN archives
\end{verbatim}

\SynapticAcknowledgments*
The project acknowledges the LaTeX3, KOMA-Script, Beamer, fontspec, and broader TeX
communities.

\end{document}
